\chapter{Variabel Acak dan Distribusi}

Sebuah layanan pesan menerima rata-rata tiga permintaan per detik pada jam sibuk.
Angka itu tercetak di dasbor, dan ia tidak menjawab pertanyaan yang mendesak: berapa
sering satu detik berisi delapan permintaan atau lebih, sehingga antrean meluap? Dua sistem bisa punya rata-rata sama persis dan berperilaku
sangat berbeda: yang satu hampir tidak pernah meluap, yang lain beberapa kali sehari.
Rata-rata saja tidak cukup, dan bukan karena hitungannya kurang teliti.

Bab 17 memberi kita bahasa untuk kejadian: peluang satu hasil, peluang gabungan,
peluang bersyarat. Bahasa itu berhenti bekerja begitu yang kita ukur berupa bilangan.
``Jumlah permintaan dalam satu detik'' bukan satu kejadian, melainkan keluarga
kejadian yang tak berhingga banyaknya, satu untuk setiap nilai yang mungkin. Untuk
cacahan, peluang tiap nilai masih bisa dijumlahkan satu per satu. Untuk besaran yang
berlanjut, cara itu runtuh: waktu tunggu bisa bernilai $0{,}317$ detik atau
$0{,}3170001$ detik, dan tidak ada cara menyebut ``peluang tepat $0{,}317$ detik''
tanpa memberi bobot nol pada setiap nilai.

Jadi dibutuhkan dua hal: cara memindahkan peluang ke garis bilangan, dan cara
mengganti penjumlahan dengan integral ketika nilainya berlanjut.
Setelah keduanya ada, rata-rata, sebaran, dan bentuk ekor bisa dihitung.

\begin{untukanak}
Bayangkan sebuah toko roti. Setiap jam, penjual mencatat berapa banyak roti yang
terjual: kadang 3, kadang 11, kadang tidak ada sama sekali. Catatan itu berupa
bilangan bulat, dan tidak bisa bernilai setengah. Kalau kita ingin tahu berapa peluang
terjual tepat 7 roti, kita bisa menghitungnya: hitung berapa jam penjualan 7 roti,
bagi dengan jumlah jam yang dicatat. Semua peluang bisa dijumlahkan satu per satu,
karena pilihannya sedikit dan bisa didaftar.

Sekarang timbang satu roti di atas timbangan digital. Hasilnya bukan bilangan bulat:
250,3 gram, atau 249,8 gram. Mau sampai berapa angka di belakang koma pun, selalu ada
kemungkinan lain yang lebih dekat lagi. Karena itu pertanyaan ``berapa peluang beratnya
tepat 250,0000 gram'' tidak berguna, sedangkan ``berapa peluang beratnya antara 249
dan 251 gram'' sangat berguna. Peluang di sini bukan milik satu titik, melainkan milik
selang.

Di bab ini kita akan tahu mengapa kedua cara itu berbeda, dan mengapa keduanya tetap
bisa memberi rata-rata yang masuk akal.
\end{untukanak}

\section{Dari hasil ke bilangan}
\tingkat{dasar}

Alat ukur tidak mencatat kejadian; ia mencatat bilangan. Jembatannya adalah fungsi
yang mengambil setiap hasil percobaan dan menetapkan satu bilangan real. Nama
``variabel acak'' menyesatkan: ia bukan variabel, dan ia tidak acak. Keacakannya
berasal dari peluang pada ruang hasil, sedangkan fungsinya biasa saja.

\begin{definisi}[Variabel acak]\index{variabel acak}
Misalkan $\Omega$ ruang hasil dan $\mathcal{F}$ kumpulan kejadian. Pemetaan
$X:\Omega\to\mathbb{R}$ disebut \emph{variabel acak} bila untuk setiap $x\in\mathbb{R}$
himpunan
\[
\{X\le x\}=\{\omega\in\Omega : X(\omega)\le x\}
\]
merupakan kejadian, yaitu anggota $\mathcal{F}$.
\end{definisi}

Syarat terakhir itu menjamin $P(X\le x)$ punya arti. Untuk hampir semua fungsi yang
muncul dalam praktik ia terpenuhi otomatis, dan hanya gagal pada penanda himpunan yang
tidak terukur.

Sebagai contoh, lempar tiga koin setimbang dan misalkan $X$ jumlah gambar yang muncul.
Nilainya $0,1,2,3$ dengan peluang $1/8,3/8,3/8,1/8$. Perhatikan bahwa $X$ memetakan
delapan hasil ke empat bilangan, jadi beberapa hasil berbagi nilai yang sama. Itulah
sebabnya peluang sebuah nilai adalah peluang gabungan kejadian.

Kita perlu satu alat yang menjawab ``sejauh mana nilainya sudah bergerak'' untuk semua
ambang sekaligus, karena ambang dipilih pembaca, bukan penulis buku.

\begin{definisi}[Fungsi distribusi kumulatif]\index{fungsi distribusi kumulatif}
Fungsi distribusi kumulatif dari variabel acak $X$ adalah
$F_X(x)=P(X\le x)$ untuk $x\in\mathbb{R}$.
\end{definisi}

\begin{proposisi}[Sifat fungsi distribusi]
Untuk setiap variabel acak $X$ berlaku:
\begin{enumerate}
  \item $F_X$ tidak turun: bila $a\le b$ maka $F_X(a)\le F_X(b)$;
  \item $\lim_{x\to-\infty}F_X(x)=0$ dan $\lim_{x\to+\infty}F_X(x)=1$;
  \item $F_X$ kontinu kanan, yaitu $F_X(x)=F_X(x^+)$ untuk setiap $x$;
  \item $P(a<X\le b)=F_X(b)-F_X(a)$ untuk setiap $a<b$.
\end{enumerate}
\end{proposisi}

\begin{proof}
Untuk (1), kejadian $\{X\le a\}$ termuat dalam $\{X\le b\}$ bila $a\le b$, dan
peluang bersifat monoton terhadap penyertaan.

Untuk (2), misalkan $A_n=\{X\le -n\}$. Kejadian ini turun, dan
$\bigcap_n A_n=\varnothing$ karena setiap hasil memberi nilai berhingga. Kekontinuan
peluang dari atas, yang menyusul dari aditif terbilang, memberi $F_X(-n)\to0$. Arah
sebaliknya memakai $B_n=\{X\le n\}$ yang naik menuju $\Omega$, sehingga
$F_X(n)\to1$.

Untuk (3), ambil barisan turun $y_n\downarrow x$ dan $C_n=\{X\le y_n\}$. Maka
$\bigcap_n C_n=\{X\le x\}$, jadi $F_X(y_n)\to F_X(x)$ menurut kekontinuan dari atas.
Karena barisan itu sebarang, limit kanan sama dengan nilai di $x$.

Untuk (4), kejadian $\{X\le b\}$ terpecah menjadi $\{X\le a\}$ dan $\{a<X\le b\}$,
sehingga $F_X(b)=F_X(a)+P(a<X\le b)$.
\end{proof}

Untuk dua lemparan koin setimbang dengan $X$ jumlah gambar,
\[
F_X(x)=\begin{cases}
0, & x<0,\\
1/4, & 0\le x<1,\\
3/4, & 1\le x<2,\\
1, & x\ge 2.
\end{cases}
\]
Fungsi ini terdefinisi di seluruh garis bilangan, walaupun $X$ hanya bernilai $0,1,2$.
Karena itu $F_X(1{,}4)=3/4$: yang ditanyakan bukan ``apakah $X$ sama dengan $1{,}4$'',
melainkan ``apakah $X$ paling besar $1{,}4$''. Setiap lompatan grafiknya adalah massa
peluang pada satu nilai. Gambar~\ref{fig:tangga-cdf} menggambar tangga itu: tiga anak
tangga mendatar, tiga lompatan tegak, dan tinggi setiap lompatan sama dengan massa
peluang di titik lompatnya.

\begin{figure}[htbp]
\centering
\begin{tikzpicture}[font=\footnotesize,>=Stealth,x=1.75cm,y=4.2cm]
  \foreach \y/\lab in {0.25/{1/4},0.75/{3/4},1/{1}}{
    \draw[layergrey!35,dashed] (-1.35,\y) -- (3.55,\y);
    \node[left,font=\scriptsize] at (0,\y) {$\lab$};
  }
  \draw[->,layergrey] (-1.35,0) -- (3.6,0) node[below] {$x$};
  \draw[->,layergrey] (0,-0.1) -- (0,1.18) node[left] {$F_X(x)$};
  \draw[very thick,layerindigo] (-1.25,0) -- (0,0);
  \draw[very thick,layerindigo] (0,0.25) -- (1,0.25);
  \draw[very thick,layerindigo] (1,0.75) -- (2,0.75);
  \draw[very thick,layerindigo] (2,1) -- (3.5,1);
  \draw[very thick,layerplum] (0,0) -- (0,0.25);
  \draw[very thick,layerplum] (1,0.25) -- (1,0.75);
  \draw[very thick,layerplum] (2,0.75) -- (2,1);
  \fill[layerindigo] (0,0.25) circle (1.8pt);
  \fill[layerindigo] (1,0.75) circle (1.8pt);
  \fill[layerindigo] (2,1) circle (1.8pt);
  \draw[layerindigo,fill=white,thick] (1,0.25) circle (1.8pt);
  \draw[layerindigo,fill=white,thick] (2,0.75) circle (1.8pt);
  \node[below,font=\scriptsize] at (0,-0.02) {$0$};
  \node[below,font=\scriptsize] at (1,-0.02) {$1$};
  \node[below,font=\scriptsize] at (2,-0.02) {$2$};
  \node[right,font=\scriptsize,text=layerplum] at (1.06,0.5) {lompatan $1/2$};
  \node[right,font=\scriptsize,text=layerplum] at (0.06,0.12) {lompatan $1/4$};
\end{tikzpicture}
\caption{Fungsi distribusi kumulatif jumlah gambar pada dua lemparan koin. Grafiknya
tangga datar di antara nilai yang mungkin, lalu melompat tepat di nilai itu. Lompatan
$1/4$ di $x=0$, $1/2$ di $x=1$, dan $1/4$ di $x=2$ adalah massa peluang $p_X(0)$,
$p_X(1)$, dan $p_X(2)$.
Titik penuh di ujung kiri setiap anak tangga menandakan $F_X$ kontinu kanan, sehingga
$F_X(1)=3/4$ dan bukan $1/4$.}
\label{fig:tangga-cdf}
\end{figure}

\kamus{variabel acak}{random variable}{fungsi yang menempelkan satu bilangan pada tiap
hasil percobaan, seperti juru tulis yang mencatat satu angka untuk setiap kejadian}

\kamus{fungsi distribusi kumulatif}{cumulative distribution function}{tangga yang
menanjak: tingginya di titik $x$ menyatakan berapa banyak peluang yang sudah
terkumpul sampai ambang itu}

\section{Fungsi massa dan fungsi kerapatan}
\tingkat{dasar}

Ada dua cara memberi bobot pada nilai variabel acak. Pada variabel yang nilainya bisa
didaftar, bobot menempel pada titik; pada yang nilainya mengisi selang, bobot menempel
pada panjang. Kekeliruan paling umum di topik ini adalah memakai cara pertama untuk
membaca hasil cara kedua.

\begin{definisi}[Variabel diskret dan fungsi massa peluang]\index{fungsi massa peluang}
$X$ disebut \emph{diskret} bila himpunan nilainya terbilang, katakanlah
$\{x_1,x_2,\dots\}$. \emph{Fungsi massa peluang} dari $X$ adalah $p_X(x)=P(X=x)$.
Fungsi ini memenuhi $p_X(x)\ge0$, $p_X(x)\le1$, dan $\sum_i p_X(x_i)=1$, serta
$F_X(x)=\sum_{x_i\le x}p_X(x_i)$.
\end{definisi}

Untuk $X$ binomial dengan $n=2$ dan $p=1/2$, massa peluangnya $1/4,1/2,1/4$ pada
$x=0,1,2$ dan nol di tempat lain. Massa peluang tidak pernah lebih besar dari satu,
karena ia memang sebuah peluang.

\begin{definisi}[Variabel kontinu dan fungsi kerapatan peluang]\index{fungsi kerapatan}
$X$ disebut \emph{kontinu} bila terdapat fungsi $f_X\ge0$ sehingga
\[
P(a\le X\le b)=\int_a^b f_X(x)\,dx
\]
untuk setiap $a\le b$, dengan $\int_{-\infty}^{\infty}f_X(x)\,dx=1$. Fungsi $f_X$
disebut \emph{fungsi kerapatan peluang}.
\end{definisi}

Kerapatan bukan peluang, dan perbedaannya bukan soal istilah. Ambil $X$ dengan
kerapatan $f_X(x)=5$ pada $[0,1/5]$ dan nol di luar itu. Maka $f_X(0{,}1)=5$, angka
yang mustahil bagi sebuah peluang, sedangkan
$P(X\le0{,}1)=\int_0^{0{,}1}5\,dx=0{,}5$. Kerapatan juga boleh tak terbatas:
$f_X(x)=\tfrac23x^{-1/3}$ pada $(0,1)$ berintegral satu walaupun meledak di dekat nol.
Satuannya berbeda: peluang per satuan panjang, seperti kepadatan penduduk adalah jiwa
per kilometer persegi. Gambar~\ref{fig:batang-dan-kurva} menaruh
kedua cara itu berdampingan: batang di kiri dijumlahkan menurut tinggi, segitiga di
kanan menurut luas, dan keduanya berjumlah satu.

\begin{figure}[htbp]
\centering
\begin{tikzpicture}[font=\footnotesize,>=Stealth]
  \begin{scope}[x=1.05cm,y=5.4cm]
    \draw[->,layergrey] (0,0) -- (5.4,0) node[right] {$x$};
    \draw[->,layergrey] (0,0) -- (0,0.62) node[left] {$p_X(x)$};
    \foreach \x/\p in {1/0.1,2/0.2,3/0.3,4/0.4}{
      \fill[layerorange!40] (\x-0.3,0) rectangle (\x+0.3,\p);
      \draw[layerbrown,thick] (\x-0.3,0) rectangle (\x+0.3,\p);
      \node[below,font=\scriptsize] at (\x,0) {$\x$};
    }
    \node[font=\scriptsize,text=layerbrown,align=left,anchor=west] at (0.45,0.44)
      {empat batang,\\ jumlah tinggi $=1$};
  \end{scope}
  \begin{scope}[xshift=7.7cm,x=1.05cm,y=5.4cm]
    \draw[->,layergrey] (0,0) -- (5.4,0) node[right] {$x$};
    \draw[->,layergrey] (0,0) -- (0,0.62) node[left] {$f_X(x)$};
    \fill[layerteal!25] (0.4,0) -- (2.4,0.5) -- (4.4,0) -- cycle;
    \draw[layerteal,thick] (0.4,0) -- (2.4,0.5) -- (4.4,0) -- cycle;
    \draw[layergrey,dashed] (2.4,0) -- (2.4,0.5);
    \node[above,font=\scriptsize] at (2.4,0.5) {$0{,}5$};
    \node[below,font=\scriptsize] at (0.4,0) {$0$};
    \node[below,font=\scriptsize] at (4.4,0) {$4$};
    \node[font=\scriptsize,text=layerteal,align=center] at (2.4,0.14)
      {luas $=\tfrac12\cdot4\cdot0{,}5=1$};
  \end{scope}
\end{tikzpicture}
\caption{Dua cara memberi bobot yang sama-sama berjumlah satu. Kiri: empat batang
setinggi $0{,}1$, $0{,}2$, $0{,}3$, $0{,}4$ pada nilai $1,2,3,4$, dijumlahkan menurut
tinggi. Kanan: kerapatan berbentuk segitiga dengan puncak $0{,}5$ pada $[0,4]$,
dijumlahkan menurut luas, dan $\tfrac12\cdot4\cdot0{,}5=1$. Tinggi $0{,}5$ di panel
kanan bukan peluang apa pun; yang berpeluang adalah luas potongan di bawahnya.}
\label{fig:batang-dan-kurva}
\end{figure}

\begin{bacakeras}
Bacakan berpasangan, satu orang memilih suara Tinggi dan satu orang suara Berat. Mulai
pelan, lalu percepat.

Tinggi: Saya tinggi kurva. Saya menjulang di satu titik, dan semua orang menunjuk saya.
Berat: Saya berat di bawah kurva. Saya menumpuk sepanjang selang, dan semua orang
menimbang saya.
Tinggi: Angka saya boleh lima. Boleh sepuluh. Boleh tak terbatas.
Berat: Angka saya tidak pernah melebihi satu, karena saya adalah seluruh peluang.
Tinggi: Tanpa saya, kau tidak tahu di mana nilai paling sering muncul.
Berat: Tanpa saya, kau tidak tahu apakah puncak itu berisi apa pun.
Bersama: Tinggi menunjuk tempat. Berat menimbang isi.

Ulangi bersama tiga kali, makin pelan: peluang itu berat, bukan tinggi.
\end{bacakeras}

\begin{teorema}\index{kerapatan}
Bila $X$ kontinu dengan kerapatan $f_X$, maka $F_X(x)=\int_{-\infty}^{x}f_X(t)\,dt$,
dan $f_X(x)=F_X'(x)$ di setiap titik $x$ tempat $F_X$ terturunkan.
\end{teorema}

\begin{proof}
Persamaan pertama adalah definisi kerapatan dengan $a=-\infty$ dan $b=x$. Persamaan
kedua adalah teorema dasar kalkulus bagian pertama: turunan fungsi akumulasi sama
dengan $f_X$ di setiap titik kekontinuan $f_X$.
\end{proof}

Pertanyaan yang paling sering menimbulkan kebingungan: bila $X$ kontinu maka
$P(X=x)=0$ untuk setiap $x$, sebab $\int_x^x f_X=0$. Apakah itu berarti kejadian $X=x$
mustahil? Tidak. Himpunan $\{\omega:X(\omega)=x\}$ tetap ada dan tetap bisa berisi
hasil; yang nol adalah peluangnya, bukan keberadaannya.

Kesulitan muncul saat mencoba menjumlahkan semua titik. Selang $[0,1]$ adalah gabungan
tak terbilang dari titik berpeluang nol, tetapi $P(0\le X\le1)=1$. Tidak ada
kontradiksi: aksioma ketiga Bab 17 hanya menjamin aditif untuk gabungan
\emph{terbilang}, dan integral mengukur panjang, bukan menghitung titik.

Ada lapisan praktis. Pembangkit bawaan NumPy membentuk nilai
sebagai kelipatan $2^{-53}$, jadi hanya ada $2^{53}$ keluaran mungkin pada $[0,1)$, dan
peluang satu nilai tertentu muncul dalam $10^7$ tarikan kira-kira
$10^7/2^{53}=1{,}1\times10^{-9}$. Pada komputer ``peluang nol'' sebenarnya
$2^{-53}\approx1{,}1\times10^{-16}$; nol yang kita tulis adalah pengidealan.

\begin{contoh}
Misalkan $X$ seragam pada $(0,1)$, dengan $f_X(x)=1$ untuk $0<x<1$. Maka
\[
P(0{,}49<X<0{,}51)=\int_{0{,}49}^{0{,}51}1\,dx=0{,}02,
\qquad P(X=0{,}5)=0 .
\]
Kerapatan di sekitar $0{,}5$ sama besar dengan di sekitar $0{,}01$; yang membedakan
hanya panjang selangnya. Karena itu pertanyaan yang benar untuk variabel kontinu
selalu berbentuk selang.
\end{contoh}

\begin{tangan}
Bahan: meteran badan, gunting, kertas karton, dan sepuluh kartu kosong.

Ukur tinggi sepuluh teman sampai sentimeter terdekat, lalu tulis satu tinggi pada tiap
kartu. Buat lima kotak di lantai dari potongan karton selebar lima sentimeter:
$150$--$154$, $155$--$159$, sampai $170$--$174$. Letakkan setiap kartu di kotak yang
cocok, bertumpuk ke atas. Yang harus kita lihat: tumpukan itu punya bentuk, dan jumlah
kartu di semua kotak tetap sepuluh.

Potong ulang karton menjadi kotak selebar sepuluh sentimeter, susun kembali kartunya,
dan lihat apa yang berubah. Tumpukan menjadi hampir separuh tingginya, tetapi tidak ada
satu kartu pun yang hilang. Lebar kotak mengubah tinggi tumpukan; jumlah kartu tidak
berubah. Itulah bedanya tinggi dan berat, dan itulah sebabnya kerapatan harus dikalikan
lebar selang sebelum boleh disebut peluang.
\end{tangan}

\begin{untukanak}
Bayangkan lapangan sepak bola yang penuh orang. Kita ingin tahu di mana orangnya
paling banyak. Kita tidak bisa menjawabnya dengan bertanya ``berapa orang tepat di
titik ini''. Jawabannya selalu nol, karena satu titik tidak punya luas. Yang bisa
dijawab adalah ``berapa orang di setiap meter persegi''. Angka itu disebut kepadatan,
dan nilainya berbeda-beda: di depan pintu masuk padat, di pojok lapangan lengang.

Kepadatan tinggi bukan berarti di situ ada banyak orang di satu titik. Kepadatan
tinggi berarti di sekitar situ, dalam wilayah kecil, orangnya menumpuk. Kalau kita
mengalikan kepadatan dengan luas wilayah, barulah keluar jumlah orang. Hal yang sama
berlaku untuk peluang: kerapatan dikali panjang selang memberi peluang, dan kerapatan
sendiri bukan peluang.

Inilah sebabnya dua kurva yang bentuknya sangat berbeda bisa memberi total yang sama,
yaitu satu. Yang penting bukan tinggi kurva di satu titik, melainkan luas di bawahnya.
Di bab ini kita akan tahu mengapa pertanyaan tentang titik harus diganti pertanyaan
tentang selang.
\end{untukanak}

\section{Harapan dan variansi}
\tingkat{menengah}

Kalau seluruh distribusi harus diringkas menjadi satu bilangan, bilangan yang dipilih
adalah titik tumpunya: letak pusat massa bila peluang dianggap berat.

\begin{definisi}[Harapan]\index{harapan}
Harapan dari variabel acak $X$ adalah
\[
\mathbb{E}[X]=\sum_i x_i\,p_X(x_i)
\qquad\text{atau}\qquad
\mathbb{E}[X]=\int_{-\infty}^{\infty}x\,f_X(x)\,dx,
\]
bergantung pada apakah $X$ diskret atau kontinu. Harapan ada hanya bila
$\mathbb{E}[|X|]<\infty$; bila tidak, dikatakan harapan $X$ tidak ada.
\end{definisi}

Syarat nilai mutlak itu bukan kehati-hatian berlebihan. Untuk distribusi Cauchy dengan
kerapatan $f(x)=1/(\pi(1+x^2))$, $\int|x|f(x)dx$ berdivergen, sehingga $\mathbb{E}[X]$
tidak ada walaupun distribusinya simetris terhadap nol; rata-rata sampelnya tidak
memusat ke mana pun, seberapa banyak pun data ditambah. Bab 19 memakai contoh ini
sebagai batas hukum bilangan besar.

Sebagai hitungan pendek, untuk $X$ Bernoulli dengan $P(X=1)=p$,
$\mathbb{E}[X]=0\cdot(1-p)+1\cdot p=p$. Untuk $X$ seragam pada $(0,1)$,
$\mathbb{E}[X]=\int_0^1x\,dx=1/2$.

Seringkali yang dibutuhkan bukan harapan $X$, melainkan harapan fungsi $X$, misalnya
$\mathbb{E}[X^2]$. Jalan lurusnya adalah mencari kerapatan $X^2$ lebih dahulu, lalu
mengintegralkan. Ada jalan yang jauh lebih pendek.

\begin{teorema}[Harapan fungsi variabel acak]
Untuk fungsi $r$ dan variabel acak $X$,
\[
\mathbb{E}[r(X)]=\sum_i r(x_i)\,p_X(x_i)
\qquad\text{atau}\qquad
\mathbb{E}[r(X)]=\int_{-\infty}^{\infty}r(x)\,f_X(x)\,dx .
\]
\end{teorema}

\begin{proof}
Untuk kasus diskret, kelompokkan nilai menurut $y=r(x)$:
\[
\sum_i r(x_i)p_X(x_i)=\sum_y y\sum_{i:\,r(x_i)=y}p_X(x_i)=\sum_y y\,P(r(X)=y),
\]
dan bentuk terakhir adalah definisi $\mathbb{E}[r(X)]$ karena $r(X)$ diskret dengan
massa peluang $P(r(X)=y)$. Untuk kasus kontinu, buktinya adalah rumus perubahan
variabel pada integral di Bab 15; pengelompokan yang sama berlaku dengan integral
menggantikan jumlah.
\end{proof}

\begin{proposisi}[Linearitas harapan]
Untuk variabel acak $X,Y$ dan konstanta $a,b$ berlaku
$\mathbb{E}[aX+bY]=a\,\mathbb{E}[X]+b\,\mathbb{E}[Y]$.
\end{proposisi}

\begin{proof}
Untuk kasus diskret dengan massa bersama $p(x,y)$,
\[
\mathbb{E}[aX+bY]=\sum_{x,y}(ax+by)p(x,y)
=a\sum_x x\sum_y p(x,y)+b\sum_y y\sum_x p(x,y)=a\,\mathbb{E}[X]+b\,\mathbb{E}[Y].
\]
Kasus kontinu identik dengan integral menggantikan jumlah.
\end{proof}

Yang tidak dipakai di sini adalah kebebasan: linearitas berlaku untuk variabel apa pun,
termasuk yang bergantung kuat. Karena itu ia aman untuk jumlah penanda, tetapi tidak
boleh dipakai untuk perkalian: $\mathbb{E}[XY]=\mathbb{E}[X]\mathbb{E}[Y]$ menuntut
kebebasan.

Harapan memberi pusat, tetapi tidak memberi lebar; ukuran lebar yang paling dipakai
adalah rata-rata kuadrat jarak ke pusat.

\begin{definisi}[Variansi]\index{variansi}
Variansi dari $X$ dengan $\mu=\mathbb{E}[X]$ adalah
$\operatorname{Var}(X)=\mathbb{E}\bigl[(X-\mu)^2\bigr]$, dan simpangan bakunya
$\sigma_X=\sqrt{\operatorname{Var}(X)}$.
\end{definisi}

\begin{proposisi}
Bila variansi ada, maka $\operatorname{Var}(X)=\mathbb{E}[X^2]-\mu^2$ dan
$\operatorname{Var}(aX+b)=a^2\operatorname{Var}(X)$.
\end{proposisi}

\begin{proof}
Kembangkan kuadratnya dan pakai linearitas:
\[
\mathbb{E}[(X-\mu)^2]=\mathbb{E}[X^2]-2\mu\,\mathbb{E}[X]+\mu^2
=\mathbb{E}[X^2]-\mu^2 .
\]
Untuk sifat kedua, $\mathbb{E}[aX+b]=a\mu+b$, sehingga
$\operatorname{Var}(aX+b)=\mathbb{E}[(aX+b-a\mu-b)^2]=a^2\mathbb{E}[(X-\mu)^2]$.
\end{proof}

Untuk $X$ Bernoulli, $\mathbb{E}[X^2]=p$ karena $X^2=X$ untuk nilai $0$ dan $1$,
sehingga $\operatorname{Var}(X)=p-p^2=p(1-p)$. Untuk $X$ seragam pada $(0,1)$,
$\mathbb{E}[X^2]=1/3$, sehingga $\operatorname{Var}(X)=1/3-1/4=1/12\approx0{,}083333$.
Variansi bersatuan kuadrat, sedangkan simpangan baku bersatuan sama dengan $X$. Satu
peringatan yang mahal: $\mathbb{E}[X^2]\ne\mathbb{E}[X]^2$, dan selisihnya justru
variansi; untuk seragam di atas, $0{,}25$ berbanding $0{,}333333$.

\begin{resep}
Bahan: tabel nilai $x_i$ dan peluang $p_X(x_i)$.
\begin{enumerate}
  \item Jumlahkan seluruh $p_X(x_i)$. Kalau hasilnya bukan satu, tabelnya belum sah.
  \item Kalikan setiap $x_i$ dengan $p_X(x_i)$, lalu jumlahkan: itulah $\mathbb{E}[X]$.
  \item Kuadratkan setiap $x_i$ lebih dahulu, kalikan dengan $p_X(x_i)$, lalu jumlahkan:
        itulah $\mathbb{E}[X^2]$.
  \item Kurangkan kuadrat hasil langkah 2 dari hasil langkah 3, yaitu
        $\operatorname{Var}(X)=\mathbb{E}[X^2]-(\mathbb{E}[X])^2$; akar kuadratnya
        adalah simpangan baku.
\end{enumerate}
Contoh: untuk $x=0,1,2,3$ dengan $p=0{,}1;\,0{,}2;\,0{,}3;\,0{,}4$, langkah 2 memberi
$2$, langkah 3 memberi $5$, dan langkah 4 memberi $1$.

Peringatan urutan: kuadratkan $x_i$ sebelum mengalikan dengan peluang, dan kurangkan
$\mu^2$ sesudah $\mathbb{E}[X^2]$ selesai dijumlahkan. Menukar urutan itu menghasilkan
angka yang tidak mengukur lebar apa pun.
\end{resep}

\kamus{harapan}{expectation}{titik tumpu penggaris berbeban: satu bilangan yang
menandai di mana nilai-nilai menumpuk}

\kamus{variansi}{variance}{ukuran seberapa jauh nilai biasa menyimpang dari titik
tumpunya, dihitung dari kuadrat jarak agar tanda plus dan minus tidak saling menghapus}

\section{Empat distribusi diskret}
\tingkat{menengah}

Empat keluarga distribusi diskret cukup untuk sebagian besar pekerjaan praktis: satu
percobaan biner, jumlah percobaan biner, waktu tunggu sampai sukses pertama, dan
kejadian langka pada selang waktu.

\paragraph{Bernoulli.} $X$ Bernoulli dengan parameter $p\in[0,1]$ bernilai $1$ dengan
peluang $p$ dan $0$ dengan peluang $1-p$, sehingga $p_X(x)=p^x(1-p)^{1-x}$ untuk
$x\in\{0,1\}$. Dari dua titik itu, $\mathbb{E}[X]=p$ dan
$\operatorname{Var}(X)=p(1-p)$; variansi terbesar pada $p=1/2$, yaitu $1/4$, tepat
ketika kedua kemungkinan seimbang.

\paragraph{Binomial.} Bila $n$ percobaan Bernoulli bebas dengan peluang sukses $p$
dijalankan, jumlah keberhasilan $X$ binomial dengan
\[
p_X(k)=\binom{n}{k}p^k(1-p)^{n-k},\qquad k=0,1,\dots,n .
\]
Rata-rata dan variansinya paling mudah diturunkan dengan memecah $X$ menjadi penanda,
bukan dari rumus itu.

\begin{teorema}\index{distribusi binomial}
Bila $X$ binomial dengan parameter $n$ dan $p$, maka
$\mathbb{E}[X]=np$ dan $\operatorname{Var}(X)=np(1-p)$.
\end{teorema}

\begin{proof}
Tulis $X=\sum_{i=1}^{n}X_i$ dengan $X_i=1$ bila percobaan ke-$i$ berhasil dan $0$ bila
gagal. Setiap $X_i$ Bernoulli dengan parameter $p$, jadi $\mathbb{E}[X_i]=p$ dan
$\operatorname{Var}(X_i)=p(1-p)$. Linearitas memberi
$\mathbb{E}[X]=\sum_i\mathbb{E}[X_i]=np$ tanpa memakai kebebasan sama sekali.

Untuk variansi, kebebasan baru diperlukan. Karena percobaannya bebas, pasangan
$X_i,X_j$ dengan $i\ne j$ juga bebas, sehingga
$\operatorname{Cov}(X_i,X_j)=p^2-p^2=0$. Maka
\[
\operatorname{Var}(X)=\sum_{i=1}^{n}\operatorname{Var}(X_i)
+2\sum_{i<j}\operatorname{Cov}(X_i,X_j)=np(1-p).
\]
\end{proof}

Rumus variansi lebih rapuh daripada rumus rata-rata: yang pertama menuntut kebebasan,
yang kedua tidak. Untuk $n=20$ dan $p=0{,}35$, $\mathbb{E}[X]=7$ dan
$\operatorname{Var}(X)=4{,}55=91/20$, dengan simpangan baku $2{,}133073$; massa di
sekitar pusatnya $p_X(7)=0{,}184401$, $p_X(10)=0{,}068614$, $p_X(12)=0{,}013564$.
Ekornya kecil tetapi bukan nol: $P(X\ge12)=0{,}019579$ dan $P(X\le4)=0{,}118197$, jadi
aturan ``tolak kotak bila ada dua belas cacat atau lebih'' menolak sekitar dua persen
kotak.

\paragraph{Geometrik.} Bila percobaan Bernoulli bebas diulang sampai sukses pertama,
banyaknya percobaan yang diperlukan adalah $X$ geometrik dengan
$p_X(k)=p(1-p)^{k-1}$ untuk $k=1,2,\dots$ Jumlah massanya satu karena deret geometri:
$\sum_{k\ge1}p(1-p)^{k-1}=p/(1-(1-p))=1$. Dengan $q=1-p$,
\[
\sum_{k\ge1}kq^{k-1}=\frac{d}{dq}\frac{1}{1-q}=\frac{1}{(1-q)^2}=\frac{1}{p^2},
\]
sehingga $\mathbb{E}[X]=p\cdot p^{-2}=1/p$ dan, dengan cara serupa,
$\operatorname{Var}(X)=(1-p)/p^2$ dari $\mathbb{E}[X^2]=(2-p)/p^2$. Untuk $p=0{,}2$:
$\mathbb{E}[X]=5$, $\operatorname{Var}(X)=20$, $P(X=5)=0{,}081920$ dan
$P(X>10)=0{,}10737418$; simulasi $10^5$ tarikan memberi rata-rata $4{,}99176$ dengan
galat baku $0{,}01414$.

\paragraph{Poisson.} Untuk $X$ Poisson dengan parameter $\lambda>0$,
\[
p_X(k)=e^{-\lambda}\frac{\lambda^k}{k!},\qquad k=0,1,2,\dots
\]
Jumlah massanya satu karena deret eksponensial:
$e^{-\lambda}\sum_k\lambda^k/k!=1$. Rata-rata dan variansinya sama, dan keduanya
keluar dari satu trik. Karena $k/k!=1/(k-1)!$,
\[
\mathbb{E}[X]=\sum_{k\ge1}k\,e^{-\lambda}\frac{\lambda^k}{k!}
=\lambda e^{-\lambda}\sum_{k\ge1}\frac{\lambda^{k-1}}{(k-1)!}=\lambda .
\]
Dengan menulis $k(k-1)/k!=1/(k-2)!$ diperoleh $\mathbb{E}[X(X-1)]=\lambda^2$, sehingga
$\operatorname{Var}(X)=\lambda$. Satu parameter mengatur pusat sekaligus lebar, dan itu
membatasi Poisson: data yang variansinya jauh melebihi rata-ratanya (disebut
\emph{overdispersi}) tidak cocok dimodelkan Poisson.

Untuk $\lambda=3$: $p_X(0)=0{,}049787$, $p_X(3)=0{,}224042$,
$P(X\ge8)=0{,}011905$. Angka terakhir menjawab pertanyaan antrean di pembuka bab:
dengan tiga permintaan per detik, sekitar $1{,}2$ persen detik berisi delapan
permintaan atau lebih.

Poisson adalah limit binomial ketika $n$ besar, $p$ kecil, dan $np\to\lambda$. Untuk
$n=1000$ dan $p=0{,}003$, binomial memberi $P(0)=0{,}049563083$ sedangkan Poisson
$e^{-3}=0{,}049787068$, selisih relatif $-4{,}5\times10^{-3}$. Kesalahannya kecil tetapi
tidak nol dan tandanya berganti menurut $k$; Poisson dipakai sebagai hampiran, bukan
kesamaan.

\section{Tiga distribusi kontinu}
\tingkat{menengah}

\paragraph{Seragam.} $X$ seragam pada $(a,b)$ bila kerapatannya $1/(b-a)$ di dalam
selang itu dan nol di luar. Fungsi distribusinya naik linear dari $0$ ke $1$, dan
\[
\mathbb{E}[X]=\frac{a+b}{2},\qquad
\operatorname{Var}(X)=\frac{(b-a)^2}{12}.
\]
Faktor $12$ muncul dari $\int_a^b(x-\tfrac{a+b}{2})^2\frac{dx}{b-a}=(b-a)^2/12$; untuk
$(0,1)$ hasilnya $1/12\approx0{,}083333$, sama dengan hitungan lewat
$\mathbb{E}[X^2]$ di bagian sebelumnya.

\paragraph{Eksponensial.} $X$ eksponensial dengan parameter $\beta>0$ memiliki
kerapatan $f_X(x)=\beta^{-1}e^{-x/\beta}$ untuk $x>0$, dengan
$F_X(x)=1-e^{-x/\beta}$, $\mathbb{E}[X]=\beta$, dan $\operatorname{Var}(X)=\beta^2$.
Parameter $\beta$ adalah rata-rata waktu tunggu. Sifat yang membedakannya dari semua
distribusi kontinu lain adalah ketiadaan memori.

\begin{proposisi}[Tanpa memori]\index{distribusi eksponensial}
Untuk $s,t\ge0$ berlaku $P(X>s+t\mid X>s)=P(X>t)$.
\end{proposisi}

\begin{proof}
$P(X>u)=e^{-u/\beta}$, sehingga
\[
P(X>s+t\mid X>s)=\frac{P(X>s+t)}{P(X>s)}
=\frac{e^{-(s+t)/\beta}}{e^{-s/\beta}}=e^{-t/\beta}=P(X>t).
\]
\end{proof}

Komponen yang sudah bertahan lima tahun tidak lebih mungkin bertahan satu tahun lagi
dibanding komponen baru: masuk akal untuk proses tanpa penuaan, tidak masuk akal untuk
lampu yang aus. Untuk $\beta=2$, $P(X>3)=0{,}22313016$ dan simulasi $10^5$ tarikan
memberi $0{,}22237$;
dengan syarat $X>2$, peluang empiris $X>5$ adalah $0{,}22028$ berbanding teori
$0{,}22313$. Eksponensial adalah kembaran kontinu geometrik: bila kedatangan mengikuti
Poisson dengan laju $\lambda$, jarak antar kedatangan eksponensial dengan
$\beta=1/\lambda$.

\paragraph{Normal.} $X$ normal dengan parameter $\mu$ dan $\sigma^2>0$ memiliki
kerapatan
\[
f_X(x)=\frac{1}{\sigma\sqrt{2\pi}}
\exp\!\left(-\frac{(x-\mu)^2}{2\sigma^2}\right),\qquad x\in\mathbb{R}.
\]
Konstanta di depan bukan pilihan: ia satu-satunya nilai yang membuat luas di bawah
kurva sama dengan satu. Bahwa luasnya satu dibuktikan dengan koordinat polar di Bab 15:
tulis $I^2$ sebagai integral ganda $e^{-(x^2+y^2)/2}$, ubah ke polar, dan dapatkan
$I^2=2\pi$.

\begin{proposisi}
Bila $X\sim N(\mu,\sigma^2)$, maka $Z=(X-\mu)/\sigma\sim N(0,1)$.
\end{proposisi}

\begin{proof}
$F_Z(z)=P(X\le\mu+\sigma z)=F_X(\mu+\sigma z)$. Turunkan terhadap $z$ dengan aturan
rantai: $f_Z(z)=\sigma f_X(\mu+\sigma z)=\varphi(z)$, yaitu kerapatan normal baku.
\end{proof}

Fungsi $\Phi(z)=P(Z\le z)$ tidak punya bentuk elementer; ia dihitung numerik atau
dibaca dari tabel. Tiga angka yang layak dihafal: $P(|Z|\le1)=0{,}682689492137$,
$P(|Z|\le2)=0{,}954499736104$, $P(|Z|\le3)=0{,}997300203937$. Untuk peluang tak
simetris, kurangkan dua nilai $\Phi$: $P(1<Z<2)=0{,}135905122$. Kuantilnya,
$\Phi^{-1}(0{,}975)=1{,}959964$, muncul di setiap interval kepercayaan 95 persen pada
Bab 20.

Normal bukan distribusi yang lebih benar daripada yang lain. Ia penting karena dua
alasan berbeda: banyak besaran alam mendekatinya, dan jumlah banyak variabel bebas
cenderung mendekatinya apa pun bentuk aslinya. Alasan kedua adalah teorema limit
pusat, dibuktikan pada Bab 19.

\kamus{distribusi normal}{normal distribution}{kurva lonceng yang simetris: hampir
semua nilai menumpuk dekat pusat, dan ekornya menipis sangat cepat}

\kamus{distribusi eksponensial}{exponential distribution}{waktu tunggu pada proses
tanpa penuaan: yang sudah lama menunggu tidak lebih dekat ke kejadian berikutnya}

\section{Distribusi bersama, kovarians, dan korelasi}
\tingkat{lanjut}

Dua distribusi terpisah tidak menyimpan informasi tentang apakah kedua besaran
cenderung naik bersama. Alat untuk itu adalah distribusi bersama.

\begin{contoh}
Misalkan $X,Y\in\{0,1\}$ dengan massa bersama $P(0,0)=0{,}1$, $P(0,1)=0{,}2$,
$P(1,0)=0{,}3$, $P(1,1)=0{,}4$. Menjumlahkan baris dan kolom memberi marginal
$P(X=1)=0{,}7$ dan $P(Y=1)=0{,}6$. Bila $X$ dan $Y$ bebas, massa pada $(1,1)$
seharusnya $0{,}42$, sedangkan yang sebenarnya $0{,}4$; jadi keduanya tidak bebas,
walaupun selisihnya kecil.
\end{contoh}

Pada kasus kontinu, distribusi bersama digambarkan oleh kerapatan bersama
$f_{X,Y}\ge0$ dengan $\iint f_{X,Y}=1$, dan marginalnya diperoleh dengan
mengintegralkan satu peubah keluar: $f_X(x)=\int f_{X,Y}(x,y)\,dy$. Kebebasan berarti
$f_{X,Y}(x,y)=f_X(x)f_Y(y)$ untuk semua $x,y$.

\begin{definisi}[Kovarians dan korelasi]\index{kovarians}\index{korelasi}
Dengan $\mu_X=\mathbb{E}[X]$ dan $\mu_Y=\mathbb{E}[Y]$,
\[
\operatorname{Cov}(X,Y)=\mathbb{E}\bigl[(X-\mu_X)(Y-\mu_Y)\bigr],
\qquad
\rho(X,Y)=\frac{\operatorname{Cov}(X,Y)}{\sigma_X\sigma_Y}.
\]
\end{definisi}

Kovarians positif berarti kedua variabel cenderung menyimpang ke arah yang sama;
negatif berarti berlawanan. Korelasi adalah versi tak bersatuan dari kovarians.

\begin{proposisi}
$\operatorname{Cov}(X,Y)=\mathbb{E}[XY]-\mathbb{E}[X]\mathbb{E}[Y]$ dan
$-1\le\rho(X,Y)\le1$.
\end{proposisi}

\begin{proof}
Kembangkan perkaliannya dan pakai linearitas. Untuk batas korelasi, pandang hasil kali
dalam $\langle U,V\rangle=\mathbb{E}[UV]$ pada peubah terpusat; ketaksamaan
Cauchy-Schwarz Bab 10 memberi
$|\langle U,V\rangle|^2\le\langle U,U\rangle\langle V,V\rangle$, yaitu
$\operatorname{Cov}(X,Y)^2\le\sigma_X^2\sigma_Y^2$.
\end{proof}

\begin{teorema}
$\operatorname{Var}(X+Y)=\operatorname{Var}(X)+\operatorname{Var}(Y)
+2\operatorname{Cov}(X,Y)$.
\end{teorema}

\begin{proof}
Tulis $X+Y-\mu_X-\mu_Y=(X-\mu_X)+(Y-\mu_Y)$, kuadratkan, lalu pakai linearitas pada
tiga suku yang muncul; suku silangnya tepat dua kali kovarians.
\end{proof}

Kembali ke tabel di atas: $\mathbb{E}[X]=7/10$, $\mathbb{E}[Y]=3/5$,
$\mathbb{E}[XY]=2/5$, sehingga $\operatorname{Cov}(X,Y)=-1/50=-0{,}02$. Dengan
$\operatorname{Var}(X)=21/100$ dan $\operatorname{Var}(Y)=6/25$, korelasinya
$-0{,}089087$ dan $\operatorname{Var}(X+Y)=0{,}41$.

Kovarians nol sering disalahartikan sebagai kebebasan, padahal arahnya hanya satu:
bebas berakibat kovarians nol, sebaliknya tidak. Ambil $X$ bernilai $-1,0,1$ dengan
peluang sama dan $Y=X^2$; maka $\mathbb{E}[X]=0$ dan
$\mathbb{E}[XY]=\mathbb{E}[X^3]=0$, sehingga korelasinya nol, padahal $Y$ ditentukan
sepenuhnya oleh $X$. Kovarians hanya menangkap ketergantungan yang berbentuk garis
lurus.

\kamus{kovarians}{covariance}{ukuran apakah dua besaran cenderung naik dan turun
bersama, dihitung dari hasil kali simpangan keduanya}

\section{Transformasi variabel acak}
\tingkat{lanjut}

Pekerjaan ini muncul terus: harga diubah ke logaritma, panjang ke luas, waktu tunggu
ke laju. Pertanyaannya selalu sama, yaitu bagaimana distribusi berubah ketika
$Y=r(X)$.

Pada kasus diskret jawabannya sederhana: massa dipindahkan dari nilai lama ke nilai
baru. Misalkan $P(X=-1)=1/4$, $P(X=0)=1/2$, $P(X=1)=1/4$, dan $Y=X^2$. Maka
$P(Y=0)=1/2$ dan $P(Y=1)=P(X=-1)+P(X=1)=1/2$; nilai $Y$ hanya dua, lebih sedikit
daripada $X$, karena $r(x)=x^2$ tidak satu-satu.

Pada kasus kontinu massanya nol di setiap titik, jadi yang dipindahkan adalah panjang,
dan panjang berubah skala. Tiga langkah baku menjawab semuanya: tentukan
$A_y=\{x:r(x)\le y\}$, hitung $F_Y(y)=\int_{A_y}f_X(x)\,dx$, lalu turunkan
$f_Y=F_Y'$.

\begin{teorema}[Transformasi monoton]\index{transformasi variabel acak}
Bila $r$ monoton tegas dan terturunkan, dengan invers $s=r^{-1}$, maka
\[
f_Y(y)=f_X\bigl(s(y)\bigr)\left|\frac{ds(y)}{dy}\right| .
\]
\end{teorema}

\begin{proof}
Misalkan $r$ naik tegas. Maka $r(x)\le y$ setara dengan $x\le s(y)$, sehingga
$F_Y(y)=F_X(s(y))$, dan aturan rantai memberi $f_Y(y)=f_X(s(y))s'(y)$ dengan
$s'(y)>0$. Bila $r$ turun tegas, ketaksamaan berbalik: $F_Y(y)=1-F_X(s(y))$, sehingga
$f_Y(y)=-f_X(s(y))s'(y)=f_X(s(y))|s'(y)|$.
\end{proof}

Faktor $|ds/dy|$ adalah peregangan lokal: bila $r$ merenggangkan skala, kerapatan
menurun, karena peluang yang sama tersebar pada selang lebih panjang. Di dimensi lebih
dari satu faktor itu menjadi determinan Jacobian, bentuk umumnya di Bab 15.
Gambar~\ref{fig:transformasi-sumbu} memperlihatkan peregangan itu pada $Y=X^2$: empat
selang sama panjang mendarat pada selang yang makin panjang ketika $y$ mendekati satu,
sementara kerapatan $Y$ meninggi di dekat nol, tempat selangnya memendek.

\begin{figure}[htbp]
\centering
\begin{tikzpicture}[font=\footnotesize,>=Stealth]
  \draw[->,layergrey] (-0.25,0) -- (7.3,0) node[right] {$x$};
  \foreach \px/\lab in {0/0,1.75/0{,}25,3.5/0{,}5,5.25/0{,}75,7/1}{
    \draw[layergrey] (\px,0.07) -- (\px,-0.07);
    \node[above,font=\scriptsize] at (\px,0.07) {$\lab$};
  }
  \node[font=\scriptsize,text=layergrey,anchor=west] at (0.15,0.55)
    {selang sama panjang};
  \draw[->,layergrey] (-0.25,-2.6) -- (7.3,-2.6) node[right] {$y=x^2$};
  \foreach \px/\lab/\anc in {0/0/east,0.44/0{,}0625/east,1.75/0{,}25/north,
                             3.94/0{,}5625/west,7/1/west}{
    \draw[layergrey] (\px,-2.53) -- (\px,-2.67);
    \node[below,font=\scriptsize,anchor=\anc] at (\px,-2.67) {$\lab$};
  }
  \foreach \px/\py in {0/0,1.75/0.44,3.5/1.75,5.25/3.94,7/7}{
    \draw[->,layerindigo] (\px,-0.05) -- (\py,-2.55);
  }
  \node[font=\scriptsize,text=layerindigo,align=center,anchor=west] at (1.0,-1.85)
    {selang memendek,\\ kerapatan meninggi};
  \node[font=\scriptsize,text=layerindigo,align=center,anchor=east] at (6.6,-1.85)
    {selang memanjang,\\ kerapatan menurun};
  \begin{scope}[xshift=8.7cm,yshift=-2.6cm]
    \draw[->,layergrey] (0,0) -- (5.4,0) node[right] {$y$};
    \draw[->,layergrey] (0,0) -- (0,3.95) node[left] {$f_Y(y)$};
    \draw[layerplum,very thick,domain=0.12:1,samples=100,smooth,variable=\t]
      plot ({5.0*\t},{1.175/sqrt(\t)});
    \draw[layerplum,dashed] (0,0) -- (0,3.8);
    \node[below,font=\scriptsize,anchor=east] at (0,0) {$0$};
    \node[below,font=\scriptsize] at (1.25,0) {$0{,}25$};
    \node[below,font=\scriptsize,anchor=west] at (5.0,0) {$1$};
    \draw[layergrey,dotted] (5.0,0) -- (5.0,1.175);
    \node[font=\scriptsize,text=layerplum,anchor=west] at (0.25,3.45) {meledak di dekat $0$};
  \end{scope}
\end{tikzpicture}
\caption{Transformasi $Y=X^2$ untuk $X$ seragam pada $(0,1)$, digambar pada dua sumbu.
Titik $x=0{,}25$ mendarat di $y=0{,}0625$, $x=0{,}5$ di $y=0{,}25$, dan $x=0{,}75$ di
$y=0{,}5625$; selang $0{,}25$ di dekat nol memendek menjadi $0{,}0625$, sedangkan selang
yang sama di dekat satu memanjang menjadi $0{,}4375$. Panel kanan adalah kerapatan
$f_Y(y)=1/(2\sqrt{y})$: nilainya $2$ di $y=0{,}0625$, $1$ di $y=0{,}25$, dan $0{,}5$ di
$y=1$. Kerapatan meninggi tepat di tempat selang memendek, sehingga luas totalnya tetap
satu.}
\label{fig:transformasi-sumbu}
\end{figure}

Bila $r$ tidak monoton, teorema di atas tidak berlaku langsung, tetapi tiga langkah
baku tetap bekerja dengan $A_y$ berupa gabungan selang.

\begin{contoh}
Misalkan $X$ seragam pada $(-1,3)$, jadi $f_X(x)=1/4$ di dalamnya, dan $Y=X^2$
bernilai di $(0,9)$. Untuk $0<y<1$ himpunan $A_y=[-\sqrt y,\sqrt y]$ seluruhnya di
dalam $(-1,3)$, jadi $F_Y(y)=\tfrac14\cdot2\sqrt y=\sqrt y/2$; untuk $1\le y<9$
himpunan itu menjadi $[-1,\sqrt y]$, jadi $F_Y(y)=(\sqrt y+1)/4$. Turunannya
\[
f_Y(y)=\begin{cases}
\dfrac{1}{4\sqrt y}, & 0<y<1,\\[4pt]
\dfrac{1}{8\sqrt y}, & 1<y<9,\\[4pt]
0, & \text{di luar itu}.
\end{cases}
\]
Kedua cabang berjumlah satu, $\tfrac12$ dan $\tfrac12$. Kerapatan $Y$ meledak di
dekat nol: $x^2$ memampatkan selang di sana.
\end{contoh}

Bila $X$ seragam pada $(0,1)$ dan $Y=X^2$, hasilnya lebih rapi: $F_Y(y)=\sqrt y$ dan
$f_Y(y)=1/(2\sqrt y)$, yang diperiksa numerik di Lab. Bila $Y=-\ln X$, maka
$F_Y(y)=1-e^{-y}$, eksponensial dengan $\beta=1$: membuang logaritma dari bilangan
seragam menghasilkan waktu tunggu eksponensial, dan kebalikannya dipakai komputer.
Bila $U$ seragam dan $F$ fungsi distribusi yang diinginkan, maka $F^{-1}(U)$
berdistribusi $F$.

\begin{tantangan}
Sebuah catatan kuliah mengusulkan kerapatan waktu tunggu
\[
f(x)=\frac{1}{2\sqrt{x}},\qquad 0<x<4,
\]
dengan alasan ``positif, turun, dan bentuknya mirip contoh di depan''. Dari kerapatan
itu penulisnya menghitung $P(X\le4)=1$ lalu menyimpulkan bahwa distribusinya sudah sah.

Periksa klaim itu. Tentukan luas total di bawah $f$, perbaiki kerapatannya dengan dua
cara sah yang berbeda, dan jelaskan mengapa perbaikannya tidak tunggal.

Petunjuk 1: yang menentukan bukan nilai $f$ di satu titik, melainkan seluruh daerah di
bawah kurvanya.

Petunjuk 2: hitung $\int_0^4\tfrac12 x^{-1/2}\,dx$; hasilnya $2$, bukan $1$. Supaya
totalnya satu, konstanta di depan harus dibagi dua menjadi $1/4$, atau selangnya
dipersempit menjadi $(0,1)$ dengan konstanta tetap.

Kalau mentok, kembali ke bagian ``Fungsi massa dan fungsi kerapatan''.
\end{tantangan}

\begin{untukanak}
Bayangkan sebuah lampu di lantai dan sebuah bola di antara lampu dan dinding.
Bayangan bola di dinding jauh lebih besar daripada bolanya, dan makin jauh bola dari
lampu, makin besar bayangannya. Tidak ada yang bertambah banyak: jumlah cat pada bola
tetap sama, hanya tersebar pada luas yang lebih besar.

Sekarang geser bola sedikit ketika ia sudah jauh dari lampu. Bayangannya bergeser
jauh. Ketika bola masih dekat lampu, geseran yang sama hanya menggeser bayangan
sedikit. Artinya, satu sentimeter di dekat lampu menjadi banyak sentimeter di dinding,
sedangkan satu sentimeter jauh dari lampu menjadi sangat banyak sentimeter.

Mengubah variabel acak bekerja seperti itu. Ada daerah yang dipampatkan dan ada
daerah yang direnggangkan, dan kerapatan hasil harus menyesuaikan diri agar totalnya
tetap satu. Kalau tidak disesuaikan, peluang akan muncul dari udara atau hilang entah
ke mana.
\end{untukanak}

\section{Contoh hitung}
\tingkat{menengah}

Sebuah lini produksi mengisi kotak berisi 20 sekering. Setiap sekering cacat dengan
peluang $0{,}35$, saling bebas, dan kotak ditolak bila berisi dua belas cacat atau
lebih.

Dengan $n=20$ dan $p=0{,}35$, aturan penolakan memberi $P(X\ge12)=0{,}019579$ dari
ekor binomial, bukan dari hampiran normal. Nilai paling mungkin, $X=7$, pun hanya
berpeluang $18{,}4$ persen, sehingga satu kotak tidak bercerita banyak.

Simulasi $10^5$ kotak dengan \texttt{default\_rng(20261024)} memberi rata-rata
$7{,}005260$ dan variansi $4{,}552918$. Galat baku rata-ratanya $0{,}006745$, jadi
simpangannya $+0{,}78$ galat baku; galat baku variansi sekitar $0{,}020348$, jadi
simpangannya $+0{,}14$. Keduanya dekat, dan itu memang diharapkan: simulasi tidak
membuktikan rumus, ia hanya gagal menyangkalnya.

Pemeriksaan kedua menyangkut total massa peluang Poisson. Jumlah
$e^{-\lambda}\lambda^k/k!$ sampai $K$ yang cukup besar harus sama dengan satu. Dihitung
dalam ruang log agar tidak meluap:

\begin{center}\small
\begin{tabular}{rrrr}
\hline
$\lambda$ & $K$ & jumlah & $1-\text{jumlah}$ \\
\hline
$0{,}5$ & 68 & $1{,}00000000000000000$ & $0$ \\
$1$ & 73 & $1{,}00000000000000000$ & $0$ \\
$3$ & 83 & $1{,}00000000000000022$ & $-2{,}2\times10^{-16}$ \\
$10$ & 107 & $1{,}00000000000000311$ & $-3{,}1\times10^{-15}$ \\
$50$ & 194 & $0{,}99999999999999778$ & $+2{,}2\times10^{-15}$ \\
$100$ & 280 & $1{,}00000000000003886$ & $-3{,}9\times10^{-14}$ \\
\hline
\end{tabular}
\end{center}

Penyimpangan terbesar muncul pada $\lambda=100$, yaitu $3{,}9\times10^{-14}$, sekitar
175 kali epsilon mesin $2{,}2\times10^{-16}$. Itu bukan tanda rumusnya salah, melainkan
tanda penjumlahan $281$ suku berukuran sangat berbeda kehilangan beberapa digit
terakhir; empat belas angka pertama masih benar.

Pemeriksaan ketiga menyangkut total kerapatan normal, karena konstanta
$1/\sqrt{2\pi}$ mudah salah ketik. Kuadratur adaptif pada $[-8,8]$ memberi
$0{,}99999999999999878$, sehingga kekurangannya $1{,}22\times10^{-15}$. Teori
meramalkan kekurangan itu tepat sebesar ekor di luar selang, yaitu
$2\Phi(-8)=1{,}24\times10^{-15}$. Kedua angka sepakat sampai dua angka penting: yang
hilang memang ekor, bukan kesalahan konstanta. Pada seluruh garis bilangan hasilnya
$0{,}99999999999999978$ dengan perkiraan galat tercetak $1{,}0\times10^{-8}$, taksiran
yang konservatif.

Pemeriksaan keempat, transformasi: kerapatan $1/(2\sqrt y)$ berintegral
$0{,}9999999999999992$, dengan sisa $8\times10^{-16}$ dari ujung $0$ tempat kerapatan
meledak. Simulasi memberi $F_Y(0{,}09)=0{,}29944$, $F_Y(0{,}25)=0{,}50003$,
$F_Y(0{,}49)=0{,}69944$, $F_Y(0{,}81)=0{,}90057$ berbanding teori $0{,}3$, $0{,}5$,
$0{,}7$, $0{,}9$. Galat baku tiap angka sekitar $0{,}00145$ pada $p=0{,}3$ dan
$0{,}00095$ pada $p=0{,}9$; keempat simpangan di bawah satu galat baku.

\kertas{Periksa dengan tangan. Hitung $\binom{20}{7}0{,}35^7 0{,}65^{13}$ dan
bandingkan dengan $0{,}184401$; jelaskan mengapa $\sum_{k=0}^{20}p_X(k)=1$ mengikuti
teorema binomial, bukan hasil pengukuran. Terakhir: rata-rata sampel $7{,}005260$
berbeda dari $7$ sebesar $0{,}00526$. Apakah itu kesalahan program? Jawab dengan
galat baku $0{,}006745$.}

\section{Lab komputasi}
\tingkat{lanjut}

Kode berikut dijalankan dengan Python 3, NumPy 2.5.2, dan SciPy 1.18.0. Semua angka
pada tabel dan paragraf sebelumnya berasal dari eksekusi ini; benihnya tetap
$20261024$, sehingga hasilnya dapat diulang.

\begin{lab}
\begin{verbatim}
import numpy as np
from math import exp, lgamma, log, sqrt, pi
from scipy import stats
from scipy.integrate import quad

rng = np.random.default_rng(20261024)
R = 100_000

# 1. binomial: rumus eksak vs Monte Carlo
n, p = 20, 0.35
mu, var = n * p, n * p * (1 - p)
sd = sqrt(var)
x = rng.binomial(n, p, size=R)
se = sd / sqrt(R)
print("Binomial n=20 p=0.35")
print(f"  teori E={mu:.6f} Var={var:.6f} sd={sd:.6f}")
print(f"  MC    rata2={x.mean():.6f} Var={x.var(ddof=1):.6f}")
print(f"  galat baku rata2={se:.6f} z={(x.mean()-mu)/se:+.4f}")
print(f"  galat baku Var  ={var*sqrt(2/(R-1)):.6f}")
print(f"  P(X>=12)={stats.binom.sf(11, n, p):.6f}")

# 2. jumlah pmf Poisson, dihitung dalam ruang log
print("\nJumlah pmf Poisson, k=0..K")
for lam in (0.5, 1.0, 3.0, 10.0, 50.0, 100.0):
    K = int(lam + 12 * sqrt(lam) + 60)
    s = sum(exp(-lam + k*log(lam) - lgamma(k+1)) for k in range(K+1))
    print(f"  lam={lam:6.1f} K={K:4d} jumlah={s:.17f}"
          f" 1-jumlah={1-s:+.2e}")

# 3. total kerapatan normal harus 1
phi = lambda t: exp(-t*t/2) / sqrt(2*pi)
print("\nKerapatan normal baku")
print(f"  phi(0)={phi(0):.15f} 1/sqrt(2pi)={1/sqrt(2*pi):.15f}")
print(f"  quad [-8,8]  ={quad(phi, -8, 8)[0]:.17f}")
print(f"  sisa={1-quad(phi, -8, 8)[0]:.2e}"
      f"  2*Phi(-8)={2*stats.norm.cdf(-8):.2e}")
v, e = quad(phi, -np.inf, np.inf)
print(f"  quad (-inf,inf)={v:.17f} galat lapor={e:.1e}")

# 4. transformasi X~Unif(0,1)
u = rng.uniform(0, 1, R)
y = u * u
print("\nY=X^2, X~Unif(0,1)")
print(f"  integral f_Y={quad(lambda t: 1/(2*sqrt(t)), 0, 1)[0]:.16f}")
for yy in (0.09, 0.25, 0.49, 0.81):
    print(f"  F_Y({yy:.2f}) MC={np.mean(y<=yy):.5f} teori={sqrt(yy):.5f}")
z = -np.log(u)
print("Y=-ln(X), harusnya Exp(1)")
print(f"  rata2={z.mean():.5f} galat baku={1/sqrt(R):.5f}"
      f" Var={z.var(ddof=1):.5f}")
print(f"  P(Y>2) MC={np.mean(z>2):.5f} teori={exp(-2):.5f}")

# 5. geometrik dan eksponensial
g = rng.geometric(0.2, R)
print("\nGeometrik p=0.2: E=5 Var=20")
print(f"  rata2={g.mean():.5f} galat baku={sqrt(20/R):.5f}"
      f" Var={g.var(ddof=1):.5f}")
w = rng.exponential(2.0, R)
print("Eksponensial beta=2: E=2 Var=4")
print(f"  rata2={w.mean():.5f} galat baku={2/sqrt(R):.5f}"
      f" Var={w.var(ddof=1):.5f}")
print(f"  P(X>3) MC={np.mean(w>3):.5f} teori={exp(-1.5):.5f}")
print(f"  P(X>5|X>2) MC={np.mean(w[w>2]>5):.5f} teori={exp(-1.5):.5f}")
\end{verbatim}
\end{lab}

Keluaran yang diperoleh, apa adanya:

\begin{verbatim}
Binomial n=20 p=0.35
  teori E=7.000000 Var=4.550000 sd=2.133073
  MC    rata2=7.005260 Var=4.552918
  galat baku rata2=0.006745 z=+0.7798
  galat baku Var  =0.020348
  P(X>=12)=0.019579

Jumlah pmf Poisson, k=0..K
  lam=   0.5 K=  68 jumlah=1.00000000000000000 1-jumlah=+0.00e+00
  lam=   1.0 K=  73 jumlah=1.00000000000000000 1-jumlah=+0.00e+00
  lam=   3.0 K=  83 jumlah=1.00000000000000022 1-jumlah=-2.22e-16
  lam=  10.0 K= 107 jumlah=1.00000000000000311 1-jumlah=-3.11e-15
  lam=  50.0 K= 194 jumlah=0.99999999999999778 1-jumlah=+2.22e-15
  lam= 100.0 K= 280 jumlah=1.00000000000003886 1-jumlah=-3.89e-14

Kerapatan normal baku
  phi(0)=0.398942280401433 1/sqrt(2pi)=0.398942280401433
  quad [-8,8]  =0.99999999999999878
  sisa=1.22e-15  2*Phi(-8)=1.24e-15
  quad (-inf,inf)=0.99999999999999978 galat lapor=1.0e-08

Y=X^2, X~Unif(0,1)
  integral f_Y=0.9999999999999992
  F_Y(0.09) MC=0.29944 teori=0.30000
  F_Y(0.25) MC=0.50003 teori=0.50000
  F_Y(0.49) MC=0.69944 teori=0.70000
  F_Y(0.81) MC=0.90057 teori=0.90000
Y=-ln(X), harusnya Exp(1)
  rata2=0.99811 galat baku=0.00316 Var=0.99623
  P(Y>2) MC=0.13478 teori=0.13534

Geometrik p=0.2: E=5 Var=20
  rata2=4.99176 galat baku=0.01414 Var=19.92571
Eksponensial beta=2: E=2 Var=4
  rata2=1.99895 galat baku=0.00632 Var=3.98400
  P(X>3) MC=0.22237 teori=0.22313
  P(X>5|X>2) MC=0.22028 teori=0.22313
\end{verbatim}

Empat hal terbaca dari keluaran itu. Pertama, setiap rata-rata sampel menyimpang kurang
dari satu galat baku: $z=+0{,}78$ untuk binomial, sedangkan $0{,}99811$, $4{,}99176$,
dan $1{,}99895$ berbanding $1$, $5$, dan $2$ untuk tiga kasus lain. Penyimpangan
sebesar itu harus muncul; kalau nol semua, justru ada yang salah.

Kedua, variansi sampel sedikit di bawah teori pada kasus berekor panjang,
$19{,}92571$ berbanding $20$, karena satu tarikan sangat besar mendominasi jumlah
kuadrat. Ketiga, galat penjumlahan Poisson tumbuh bersama $\lambda$, dari nol menjadi
$10^{-14}$ ketika jumlah sukunya bertambah dari $74$ menjadi $281$. Keempat, kerapatan
hasil berintegral satu sampai $16$ angka, dan $F_Y$ cocok dengan $\sqrt y$ di empat
titik uji; yang kedua lebih penting, karena integral satu hanya memeriksa konstanta,
sedangkan kecocokan di empat titik memeriksa bentuk fungsi.

Kode itu tidak membuktikan satu pun teorema di bab ini; ia memeriksa bahwa rumus yang
ditulis memberi angka yang disebutkan. Bukti $\mathbb{E}[X]=np$ ada pada argumen
penanda, bukti $f_Y(y)=1/(2\sqrt y)$ pada perhitungan fungsi distribusi; keluaran
program adalah data, bukan argumen.

\begin{untukdosen}
Urutan yang saya sarankan: dua sistem dengan rata-rata sama tetapi ekor berbeda, lalu
definisi variabel acak, lalu langsung fungsi distribusi kumulatif. Kerjakan satu tabel
kecil di papan, gambar tangganya, tanyakan $F(1{,}4)$, baru setelah itu kerapatan.
Tahan diri untuk tidak menyebut integral sebelum mahasiswa melihat mengapa penjumlahan
tidak lagi cukup.

Dua miskonsepsi yang paling sering muncul. Pertama, $f_X(x)$ dibaca sebagai peluang,
sehingga kerapatan bernilai $5$ terasa mustahil dan $P(X=x)$ dianggap sama dengan
$f_X(x)$. Kedua, $P(X=x)=0$ disimpulkan berarti mustahil, sehingga
$P(0\le X\le1)=1$ menjadi paradoks. Jawabannya aditif terbilang: gabungan tak
terbilang tidak dijumlah.

Tanda siap lanjut: mahasiswa dapat menjelaskan perbedaan massa dan kerapatan tanpa
membaca catatan, dan menghitung $\mathbb{E}[X^2]$ tanpa mencari kerapatan $X^2$.
\end{untukdosen}

\section{Ringkasan}
\tingkat{dasar}

\begin{ringkasan}
\begin{itemize}
\item Variabel acak adalah fungsi dari ruang hasil ke bilangan real, bukan bilangan
      acak. Fungsi distribusi $F_X(x)=P(X\le x)$ memuat seluruh informasi distribusi:
      tidak turun, kontinu kanan, berlimit $0$ dan $1$.
\item Fungsi massa membobot titik dan nilainya paling besar satu; fungsi kerapatan
      membobot panjang, boleh melebihi satu, dan $P(X=x)=0$ di setiap titik. Peluang
      nol bukan kemustahilan, karena yang menjumlah hanya gabungan terbilang.
\item Harapan adalah titik tumpu dan linear tanpa syarat kebebasan;
      $\operatorname{Var}(X)=\mathbb{E}[X^2]-\mathbb{E}[X]^2$ adalah rata-rata kuadrat
      simpangan.
\item Binomial $\mathbb{E}=np$, $\operatorname{Var}=np(1-p)$, diturunkan lewat
      penanda; geometrik $\mathbb{E}=1/p$, $\operatorname{Var}=(1-p)/p^2$; Poisson
      $\mathbb{E}=\operatorname{Var}=\lambda$, limit binomial untuk kejadian langka.
\item Normal distandardisasi menjadi $N(0,1)$ dan dibaca lewat $\Phi$; eksponensial
      tanpa memori, kembaran kontinu geometrik.
\item Kovarians mengukur ketergantungan linear saja: nol tidak berarti bebas, dan
      transformasi monoton mengubah kerapatan dengan faktor $|ds/dy|$.
\end{itemize}
\end{ringkasan}

\section{Latihan}
\tingkat{menengah}

\begin{latihan}
\begin{enumerate}
  \item Misalkan $p_X(-2)=0{,}1$, $p_X(0)=0{,}4$, $p_X(3)=c$. Tentukan $c$, lalu
        hitung $\mathbb{E}[X]$, $\mathbb{E}[X^2]$, dan $\operatorname{Var}(X)$.
  \item Buktikan $P(a<X\le b)=F_X(b)-F_X(a)$, lalu jelaskan mengapa
        $P(a\le X\le b)$ bisa berbeda dari $P(a<X\le b)$ untuk variabel diskret dan
        mengapa keduanya sama untuk variabel kontinu.
  \item Nilai klaim ini: ``karena $P(X=x)=0$ untuk setiap $x$, kejadian $X=x$
        mustahil.'' Sebutkan bagian yang benar, bagian yang keliru, dan peran aditif
        terbilang.
  \item Turunkan $\mathbb{E}[X^2]=(2-p)/p^2$ untuk geometrik dengan menurunkan
        $\sum k^2q^{k-1}$ dua kali, lalu simpulkan $\operatorname{Var}(X)$ dan periksa
        untuk $p=0{,}2$.
  \item Misalkan $X\sim N(3,4)$. Hitung $P(X>1)$, $P(1<X<5)$, dan kuantil 0,2 untuk
        $X$; bandingkan dengan \texttt{scipy.stats.norm}.
  \item Hitung $\sum_{k=0}^{40}e^{-2{,}5}2{,}5^k/k!$ dan $P(X\ge8)$ untuk Poisson
        $\lambda=2{,}5$, lalu bandingkan dengan binomial $n=1000$, $p=0{,}0025$;
        jelaskan arah selisihnya.
  \item Buktikan eksponensial tanpa memori, lalu jelaskan mengapa geometrik juga tanpa
        memori dan mengapa seragam pada bilangan bulat tidak.
  \item Untuk $X$ seragam pada $(-1,3)$ dan $Y=X^2$, turunkan $f_Y$, periksa
        $\int f_Y=1$, lalu hitung $\mathbb{E}[Y]$ lewat $f_Y$ dan lewat
        $\mathbb{E}[X^2]$.
  \item Konstruksikan dua variabel yang kovariansnya nol tetapi tidak bebas. Hitung
        $\operatorname{Cov}$ dan $\rho$ secara eksak, dan tunjukkan peluang bersama
        yang membuktikan ketidakbebasan.
  \item Berapa ulangan $R$ agar galat baku rata-rata binomial $n=20$, $p=0{,}35$
        paling banyak $0{,}005$? Berapa pula agar galat baku taksiran variansi paling
        banyak $0{,}02$?
\end{enumerate}

\noindent\textbf{Petunjuk.} Butir 3: pisahkan pernyataan tentang peluang dari
pernyataan tentang keberadaan himpunan. Butir 4: turunkan $\sum_{k\ge0}q^k=1/(1-q)$
dua kali terhadap $q$, dan ingat $\sum k^2q^{k-1}$ memuat $\sum k(k-1)q^{k-1}$.
Butir 8: untuk $0<y<1$ himpunan $A_y$ memuat dua cabang, untuk $y\ge1$ hanya satu.
\end{latihan}

\begin{untukpeneliti}
Yang benar-benar terbukti di bab ini: sifat fungsi distribusi beserta buktinya,
hubungan kerapatan dengan fungsi distribusi lewat teorema dasar kalkulus, rumus
harapan fungsi variabel acak untuk kasus diskret, linearitas harapan, rumus variansi,
rata-rata dan variansi binomial lewat penanda, rata-rata geometrik dan Poisson lewat
deret, ketiadaan memori eksponensial, dan rumus transformasi monoton.

Yang dikutip, bukan dibuktikan: bahwa fungsi yang memenuhi tiga sifat fungsi distribusi
selalu berasal dari suatu ruang peluang (Wasserman, Teorema 2.8, yang arahnya memakai
teorema Caratheodory), bahwa fungsi distribusi menentukan seluruh peluang (Wasserman,
Teorema 2.7), normalisasi normal lewat koordinat polar (Bab 15), dan batas korelasi
lewat Cauchy-Schwarz (Bab 10).

Sumber primer: Wasserman, \emph{All of Statistics}, Bagian 2.1--2.4 (termasuk Lemma
2.15 dan peringatan bahwa kerapatan boleh melebihi satu) dan 2.11 (Contoh 2.47 untuk
$Y=X^2$ pada $U(-1,3)$), serta 3.1--3.4. Blitzstein dan Hwang, \emph{Introduction to
Probability}, Bab 3--5 (LOTUS di 4.5, Poisson di 4.7, eksponensial di 5.5) dan 7.3
untuk kovariansi.

Batas keberlakuan: kerapatan tidak tunggal, sehingga setiap pernyataan tentang $f_X$
adalah pernyataan tentang kelas fungsi yang sama hampir di mana-mana; rumus Jacobian
menuntut monoton dan terturunkan; harapan Cauchy tidak ada, sehingga setiap
``rata-rata'' pada data berekor sangat tebal harus dicurigai lebih dahulu. Pertanyaan
terbuka: seberapa baik kerapatan dapat ditaksir dari data tanpa asumsi bentuk, dan
apakah ekor distribusi nyata cukup tipis untuk membenarkan pemakaian variansi.
\end{untukpeneliti}

\peta{Bab 18 pada Wasserman, \emph{All of Statistics}, Bab 2 (2.1--2.4, Lemma 2.15,
dan 2.11 untuk transformasi) dan Bab 3 (3.1--3.4). Pada Blitzstein dan Hwang,
\emph{Introduction to Probability}, Bab 3 (3.6--3.7), Bab 4 (4.2, 4.5, 4.7), Bab 5
(5.1--5.5) dan 7.3. Integral yang dipakai adalah integral Riemann Bab 13; perubahan
variabel dimensi lebih dari satu ada di Bab 15.}
